\section*{Chapter \ref{ch:b2:climate-variability} --- Climate Variability}

\begin{solution}{exo:b2:climate-variability:1}
(a) $29.27 - 26.71 = \qty{2.56}{\celsius}$ in 2026; $28.79 - 26.73 = \qty{2.06}{\celsius}$ in 1997.
(b) September 2026 was both the warmer month, by \qty{0.48}{\celsius}, and the larger anomaly, by
\qty{0.50}{\celsius}. (c) Each year is compared with a 30-year base period centred on it, moved
every five years; the September \omterm{def:b2:climate-variability:modes}{climatology} of the 1980s and 1990s differs slightly from that of
the 2010s and 2020s.
\end{solution}

\begin{solution}{exo:b2:climate-variability:2}
(a) JJA: $(0.79 + 0.95 + 1.28)/3 = 1.007$; JAS: $(0.95 + 1.28 + 1.55)/3 = 1.260$; ASO: $(1.28 +
1.55 + 1.68)/3 = 1.503$; SON: $(1.55 + 1.68 + 2.01)/3 = 1.747$; OND: $(1.68 + 2.01 + 2.02)/3 =
1.903$, all in \unit{\celsius}. (b) Rounded to one decimal: 1.0, 1.3, 1.5, 1.7, 1.9, the values of
the table. (c) In ASO 2023, at 1.5.
\end{solution}

\begin{solution}{exo:b2:climate-variability:3}
(a) The index follows the departure of the pressure at Tahiti minus that at Darwin: below average
at Tahiti and above at Darwin, it is negative. (b) El~Niño. (c) Weaker easterly trade winds, even
westerly winds near the surface in the west, and less rain than usual over Indonesia, since the
deep convection has moved east.
\end{solution}

\begin{solution}{exo:b2:climate-variability:4}
(a) $804.1/880.6 = \qty{91.3}{\percent}$ of the average: within 90 to \qty{110}{\percent}, a
normal season, though in its lower part. (b) The season totals $556 + 659 + 427 + 252 =
\qty{1894}{mm}$; July brings $659/1894 = \qty{34.8}{\percent}$ of it, more than a third.
\end{solution}

\begin{solution}{exo:b2:climate-variability:5}
(a) $\tau = 1025 \times 3992 \times 50/16.1 = \qty{1.271e7}{s} = \qty{147.1}{\day}$ for
\qty{50}{m}, and twice that, \qty{294.2}{\day}, for \qty{100}{m}. (b) $\mathrm{e}^{-30.44/147.1} =
0.813$ and $\mathrm{e}^{-30.44/294.2} = 0.902$. (c) The damping time is proportional to $h$: a
shallower mixed layer stores less heat for the same flux and forgets an anomaly faster, so the
ocean's memory, and the predictability it gives, decline.
\end{solution}

\begin{solution}{exo:b2:climate-variability:6}
(a) Three months are 91.3 days: $1.2 \times \mathrm{e}^{-91.3/147.1} = \qty{0.65}{K}$; six months
are 182.6 days: $1.2 \times \mathrm{e}^{-182.6/147.1} = \qty{0.35}{K}$. (b) The anomaly decays
slowly, so it is a better guess of next month's anomaly than zero. But the random forcing adds
new anomalies every day, unknown at the time of the forecast; as the initial anomaly decays,
they take an ever larger share, and the forecast drifts towards the \omterm{def:b2:climate-variability:modes}{climatology}, where it no
longer helps.
\end{solution}

\begin{solution}{exo:b2:climate-variability:7}
(a) $160 \times 86\,400 = \qty{1.38e7}{J}$, \qty{13.8}{MJ} per square metre and per day. (b) The
layer holds $(1000 - 850) \times 100/9.81 = \qty{1529}{kg}$ of air per square metre, with a heat
capacity of $1529 \times 1004.5 = \qty{1.536e6}{J.m^{-2}.K^{-1}}$; it would warm by $1.38 \times
10^7/1.536 \times 10^6 = \qty{9.0}{K}$ a day. (c) The layer radiates heat in the infrared, to
space and to the air above; and the warm air rises and mixes with the layer above it, or is
carried away by the winds, among them the inflow of cooler air the heat low itself draws in.
\end{solution}

\begin{solution}{exo:b2:climate-variability:8}
(a) $(-1.93 - 1.11 - 1.98)/3 = -1.673$, about $-1.67$. (b) It is the lowest winter of the record.
(c) A strongly negative NAO: a weak Icelandic Low and Azores High, a weaker jet and a storm track
further south than usual; a cold, dry winter in northern Europe, a wet one in southern Europe,
and more outbreaks of cold air in the eastern United States.
\end{solution}

\begin{solution}{exo:b2:climate-variability:9}
(a) $9.3 - 0.4 = \qty{8.9}{mm}$ a day, more than 2.5: in the domain. (b) In the Southern Hemisphere
the summer is December--February: $6.1 - 4.2 = \qty{1.9}{mm}$ a day, less than 2.5: outside it,
although it rains more there than at most places of the domain. (c) $1894/122 = \qty{15.5}{mm}$ a
day; Kerala would be outside only if its winter rate exceeded $15.5 - 2.5 = \qty{13.0}{mm}$ a day.
(d) A \omterm{def:b2:climate-variability:monsoon}{monsoon} is a seasonal reversal: a wet place that is wet all year round, like the second
station, has no \omterm{def:b2:climate-variability:monsoon}{monsoon}, and a dry place with a short, intense rainy season has one.
\end{solution}

\begin{solution}{exo:b2:climate-variability:10}
(a) Differentiating, $\ddot T = a\dot H = -abT$: the equation of a harmonic oscillator of angular
frequency $\omega = \sqrt{ab}$, and period $2\pi/\sqrt{ab}$. (b) $2\pi/\sqrt{ab} = 4$ gives $ab =
(2\pi/4)^2 = \qty{2.47}{yr^{-2}}$. (c) With $T = u\,\mathrm{e}^{-rt}$, $H = v\,\mathrm{e}^{-rt}$,
$\dot T = (\dot u - ru)\,\mathrm{e}^{-rt} = (av - ru)\,\mathrm{e}^{-rt}$, so $\dot u = av$, and
likewise $\dot v = -bu$: $(u, v)$ oscillates without damping, and $(T, H)$ is that oscillation
multiplied by $\mathrm{e}^{-rt}$. In one period the amplitude falls by $\mathrm{e}^{-0.25 \times
4} = \mathrm{e}^{-1} = 0.37$. (d) A forcing: the random fluctuations of the trade winds, which
kick the system again and again, and the \omterm{def:b2:climate-variability:bjerknes}{Bjerknes feedback}, which makes the damping weaker or
even negative.
\end{solution}

\begin{solution}{exo:b2:climate-variability:11}
(a) $A = R^2 \,\Delta\lambda\,(\sin\ang{5} - \sin(\ang{-5})) = (6.371 \times 10^6)^2 \times 0.8727
\times 0.1743 = \qty{6.17e12}{m^2}$. (b) $Q = 1025 \times 3992 \times 100 \times 2 \times 6.17
\times 10^{12} = \qty{5.05e21}{J}$. (c) Six months are $182.6 \times 86\,400 = \qty{1.578e7}{s}$;
per square metre, $8.18 \times 10^8/1.578 \times 10^7 = \qty{51.9}{W/m^2}$, sustained for half a
year. (d) From the ocean itself: weaker trades bring up less cold water in the east, and the warm
water of the west, stored above the deeper thermocline, spreads east; much of the warming
can come from heat already in the ocean, moved rather than gained.
\end{solution}

\begin{solution}{exo:b2:climate-variability:12}
(a) The loop returns $0.3g$, then $0.3g^2$, and so on: the total is $0.3(1 + g + g^2 + \dots) =
0.3/(1 - g)$, a geometric series that converges for $g < 1$: \qty{0.6}{K}, \qty{1.5}{K} and
\qty{3.0}{K} for $g = 0.5$, 0.8 and 0.9. (b) The series diverges: the anomaly would grow without
limit; the growth must be stopped by processes the loop leaves out, the damping by the surface
heat flux and the draining of the warm-water store that the recharge oscillator describes. (c)
The warmer the sea surface, the more readily a small warming triggers deep convection and rain,
whose heating drives the winds; that is the season in which the effects of ENSO are strongest.
\end{solution}

\begin{solution}{pb:b2:climate-variability:1}
\textbf{1.} $(0.26 + 0.75 + 1.13)/3 = \qty{0.713}{\celsius}$.
\textbf{2.} MJJ $(0.75 + 1.13 + 1.52)/3 = 1.133$; JJA $(1.13 + 1.52 + 1.80)/3 = 1.483$; JAS $(1.52 +
1.80 + 2.06)/3 = 1.793$.
\textbf{3.} ASO $(1.80 + 2.06 + 2.26)/3 = 2.040$; SON $(2.06 + 2.26 + 2.42)/3 = 2.247$; OND $(2.26
+ 2.42 + 2.36)/3 = 2.347$; NDJ $(2.42 + 2.36 + 2.33)/3 = 2.370$.
\textbf{4.} Rounded to one decimal, 0.7, 1.1, 1.5, 1.8, 2.0, 2.2, 2.3 and 2.4: the table.
\textbf{5.} A single month carries weather noise; ENSO lasts many months, and the mean of three
removes much of the noise while keeping the slow signal.
\textbf{6.} In JJA 1997, at 1.483, rounded to 1.5; it peaked at 2.4 in NDJ.
\textbf{7.} From ASO 2018 to MJJ 2019, and from SON 2019 to FMA 2020.
\textbf{8.} Ten seasons, and six.
\textbf{9.} In ASO 2020; within 2020 only four seasons, ASO to NDJ, are at or below $-0.5$, one
short of the rule, and DJF 2021 completes it.
\textbf{10.} One El~Niño, SON 2018 to MAM 2019, seven seasons (AMJ 2019 falls to 0.37); the
La~Niña starts in AMJ 2020 ($-0.55$), four seasons before the ONI places it, in ASO.
\textbf{11.} No: the relative index never exceeds 0.29 in late 2019. The whole tropical ocean was
then warmer than its base period by about half a degree; the Niño-3.4 region was warm in
absolute terms, but hardly warmer than the tropics around it.
\textbf{12.} Eleven.
\textbf{13.} 8, 7, 10, 4, 5, 6, 12, 5, 9 and 3 years.
\textbf{14.} The shortest is 3 years (2023 to 2026), the longest 12 (1997 to 2009).
\textbf{15.} $69/10 = 6.9$ years either way: the sum of successive intervals is the time from the
first start to the last, $2026 - 1957 = 69$ years.
\textbf{16.} It has already lasted five seasons, MAM to JAS, at or above \qty{0.5}{\celsius}, and
peaked at 2.2, above 1.5; its start, MAM 2026, is the first season of the run, and later seasons
can only extend the run, not move its beginning.
\textbf{17.} No, and it need not be: the 25 episodes of all strengths come about every three years, as the FAQ says;
the strong ones are a subset, so they are rarer, and their intervals range from 3 to 12 years.
\textbf{18.} $\tau = 1025 \times 3992 \times 50/16.1 = \qty{1.27e7}{s} = \qty{147}{\day}$.
\textbf{19.} Mixed layer: $1025 \times 3992 \times 50 = \qty{2.05e8}{J.m^{-2}.K^{-1}}$; air column:
$(10^5/9.81) \times 1004.5 = 10\,194 \times 1004.5 = \qty{1.02e7}{J.m^{-2}.K^{-1}}$, 20 times
less.
\textbf{20.} $1.024 \times 10^7/(1025 \times 3992) = \qty{2.5}{m}$ of seawater.
\textbf{21.} $1.024 \times 10^7/16.1 = \qty{6.36e5}{s} = 7.4$~days.
\textbf{22.} An anomaly of the air alone is damped in about a week, one of the mixed layer in
about five months, and the warm layer above the thermocline is deeper still: only the ocean holds
heat long enough to carry an anomaly from one season to the next, while the atmosphere supplies
the winds that couple the two.
\textbf{23.} Since 1950, strong El~Niño episodes have begun on average every \textbf{6.9 years},
about seven.
\end{solution}
